%==============================================
% VOLUME III: Array Odyssey
%==============================================
\part{Array Odyssey}
\label{part:array-odyssey}

\begin{volumesintro}
\lettrine[lines=3]{A}{t last} we arrive at arrays themselves---the most fundamental data structure in computing, and the destination of everything Volumes I and II built up to. This volume explores arrays from theory to implementation to the structures and algorithms built on top of them, the parallel and distributed systems that process them at scale, and a representative tour---not an encyclopedia---of the domains where arrays do real work, from machine learning to scientific computing.

\vspace{1em}
\textbf{What Makes This Different:}
\begin{itemize}[noitemsep]
    \item \textbf{Focused Coverage:} Every major technique and structure built on arrays, without the redundant variants that bloat a reference into an encyclopedia
    \item \textbf{Hardware Awareness:} Cache effects, alignment, SIMD, and GPU execution---grounded in Volume II
    \item \textbf{Rigorous Analysis:} Mathematical foundations and performance, from a single access to a distributed cluster
    \item \textbf{Practical Optimization:} Making arrays fast in real systems, from one core to many
\end{itemize}

\begin{quote}
\textit{``Simplicity is prerequisite for reliability.''}

\hfill--- \textsc{Edsger W. Dijkstra}
\end{quote}
\end{volumesintro}

% ============================================================
% STAGE 1: Array Theory, Memory Layout, and Lifecycle
%
%   Goal check: after this stage, the reader can explain what an array
%   formally is, how it sits in memory, and how its lifecycle (alloc,
%   init, copy, grow, shrink) is managed -- the vocabulary every later
%   stage assumes.
% ============================================================

\chapter{Mathematical Foundations of Arrays}
Formal definition, array as function from index set to value set, properties, mathematical operations.

\chapter{Array as Abstract Data Type}
Interface specification, invariants, preconditions, postconditions, ADT vs. implementation.

\chapter{One-Dimensional Array: Memory Layout}
Contiguous allocation, base address, element offset calculation, memory alignment.

\chapter{Address Calculation and Pointer Arithmetic}
Computing element addresses, stride, pointer arithmetic rules, bounds.

\chapter{Zero-Based vs. One-Based Indexing}
Historical reasons, language conventions, conversion, pros and cons.

\chapter{Array Bounds and Boundary Conditions}
Index range, out-of-bounds access, undefined behavior, bounds checking strategies.

\chapter{Array Allocation: Stack vs. Heap}
Automatic (stack) allocation, dynamic (heap) allocation, lifetime management, size limitations.

\chapter{Static Array Declaration}
Compile-time size, storage duration, initialization, usage patterns.

\chapter{Dynamic Memory Allocation and Deallocation}
malloc/calloc/realloc, new/delete, proper cleanup (RAII, smart pointers), memory leaks and fragmentation.

\chapter{Array Initialization: Techniques and Performance}
Zero-initialization, value and aggregate initialization, designated initializers; memset internals, compiler optimizations, lazy initialization costs.

\chapter{Array Copying and Memory Transfer}
Shallow vs. deep copy, copy semantics and assignment operators; memcpy/memmove, bulk copy performance, vectorized copying.

\chapter{Array Access Patterns and Cache Line Utilization}
Sequential, random, and strided access; impact on cache and prefetching; cache line size, spatial locality, utilizing full cache lines.

\chapter{Prefetching and Traversal Optimization}
Hardware and software prefetching, prefetch distance; loop order, blocking/tiling, minimizing cache misses during traversal.

\chapter{Dynamic Arrays: Implementation and Growth Strategies}
Capacity vs. size, C++ \texttt{std::vector} internals; doubling, multiplicative, and additive growth, and their tradeoffs. (Amortized-cost proof techniques for this growth are developed properly in Stage 5.)

\chapter{Dynamic Array Reallocation and Shrinking}
In-place vs. move-to-new-location reallocation, minimizing copying; shrink-to-fit, hysteresis, avoiding oscillation.

\chapter{Dynamic Arrays in Practice}
C++ \texttt{std::vector}, Java \texttt{ArrayList}, Python \texttt{list} compared; iterator invalidation on reallocation and safe iteration patterns.

\chapter{Multidimensional Arrays: Layout Conventions}
Dimensions and shape; row-major (C) vs. column-major (Fortran/BLAS/LAPACK) order; linear index from multiple indices, stride computation.

\chapter{Multidimensional Array Traversal}
Loop nesting order, cache-friendly traversal, impact of memory layout on iteration performance.

\chapter{Array of Arrays, Jagged Arrays, and Dope Vectors}
Pointer indirection vs. true multidimensional storage; non-rectangular (jagged) arrays; dope vectors and array descriptors as metadata for flexible array operations.

\chapter{Sparse Arrays: Concepts and Construction-Friendly Formats}
When most elements are zero; coordinate list (COO), dictionary of keys (DOK), and list of lists (LIL) as formats optimized for easy construction.

\chapter{Sparse Arrays: Compressed Formats and Operations}
Compressed sparse row (CSR) and column (CSC) for efficient row/column operations; sparse matrix-vector and sparse-sparse operations; the fill-in problem.

\chapter{Bit Arrays and Bitmap Indexes}
Packing bits into words; computing word/bit offsets; bitwise operations on arrays; bitmaps for indexing, rank and select operations.

\chapter{Circular Arrays and Ring Buffers}
Wrap-around indexing, modulo arithmetic; tracking fill level, full vs. empty distinction, producer-consumer patterns.

% ============================================================
% STAGE 2: Searching and Sorting
%
%   Goal check: after this stage, the reader can choose and implement the
%   right search or sort for a given constraint (memory, stability,
%   data distribution, I/O-boundedness) rather than defaulting to
%   whatever they memorized.
% ============================================================

\chapter{Linear and Binary Search}
Sequential scan with early termination and the sentinel technique; binary search's prerequisites, algorithm, loop invariants, and complexity analysis.

\chapter{Binary Search: Variants and Applications}
Iterative vs. recursive, overflow-safe midpoint, lower\_bound vs. upper\_bound; finding insertion points, counting occurrences, searching rotated arrays.

\chapter{Specialized Search Algorithms}
Interpolation search (uniform distributions), exponential search (unbounded ranges), ternary search (unimodal functions), jump search, and Fibonacci search, compared on their assumptions and when each beats binary search.

\chapter{Sorting Fundamentals}
Comparison-based lower bound, stability, adaptivity, in-place vs. out-of-place sorting.

\chapter{Elementary Sorts: Bubble, Selection, Insertion, and Shell}
Adjacent swaps and early termination; selection by repeated minimum; insertion into a sorted prefix (adaptive, online); Shell sort's gap sequences and diminishing increments.

\chapter{Merge Sort: Algorithm and Variants}
Divide-and-conquer with guaranteed $O(n \log n)$ and stability; top-down vs. bottom-up implementation; insertion-sort cutover and copy-elimination optimizations.

\chapter{External Merge Sort}
Sorting data larger than memory; multi-way merge, minimizing I/O, the connection to B-trees.

\chapter{Quick Sort: Algorithm, Pivots, and Partitioning}
Partitioning and divide-and-conquer; pivot selection strategies (first, random, median-of-three, ninther); Lomuto, Hoare, and three-way partitioning; tail-recursion and small-subarray optimizations.

\chapter{Randomized Quicksort and Introsort}
Random pivot selection and expected-case analysis; the hybrid introsort strategy (quicksort with depth limiting and a heapsort fallback) used in \texttt{std::sort}.

\chapter{Heap Sort and the Binary Heap}
Building a heap and repeated extract-max; the complete binary tree stored in an array, parent/child index formulas, sift-up/sift-down, and bottom-up heap construction in $O(n)$.

\chapter{Counting Sort and Radix Sort}
Counting sort's linear time via cumulative counts; radix sort's digit-by-digit approach (LSD vs. MSD), and why the per-digit sub-sort must be stable.

\chapter{Bucket Sort}
Distributing into buckets under a uniform-distribution assumption; expected $O(n)$ performance.

\chapter{Timsort}
Python's adaptive merge sort: detecting and merging existing runs, galloping mode.

\chapter{pdqsort (Pattern-Defeating Quicksort)}
Rust's default sort: detecting adversarial and already-sorted patterns, block partitioning, fallback strategies.

\chapter{Sorting Networks}
Comparator networks, Batcher's sort, bitonic sort, and why fixed, data-independent comparison sequences matter for hardware and parallel sorting.

% ============================================================
% STAGE 3: In-Place Manipulation, Two Pointers, and Prefix Sums
%
%   Goal check: after this stage, the reader recognizes a wide class of
%   interview-style and real-world array problems as instances of a small
%   number of reusable patterns, instead of memorizing each one separately.
% ============================================================

\chapter{Array Rotation: Algorithms and Tradeoffs}
Reversal-based rotation (three reversals, elegant proof), the juggling algorithm (GCD-based cycles), and block-swap rotation; in-place reversal via the two-pointer technique.

\chapter{Permutations: Application, Generation, and Shuffling}
Applying a permutation in place via cycle decomposition; generating all permutations (Heap's algorithm, lexicographic order); the Fisher-Yates shuffle and why naive shuffles are biased.

\chapter{Partitioning Algorithms}
Two-way (quicksort-style), three-way (Dutch flag), and $k$-way partitioning; stable partitioning and its implementation strategies.

\chapter{Selection Algorithms: Quickselect and Median of Medians}
Finding the $k$-th smallest element via partition-based selection (expected $O(n)$); the deterministic linear-time median-of-medians construction and why it matters theoretically.

\chapter{Two-Pointer Technique and Sliding Windows}
Meeting-in-the-middle and opposite-direction pointers; fixed- and variable-size sliding windows, and maintaining window state (including deque-based min/max).

\chapter{Multi-Pointer Techniques}
Three-pointer and beyond: the Dutch National Flag as a three-pointer problem, merging sorted arrays, and complex partitioning that two pointers can't express.

\chapter{Prefix Sums: Construction and Applications}
Building the cumulative-sum array in $O(n)$ for $O(1)$ range-sum queries; applications to subarray sum and equilibrium-index problems.

\chapter{Two-Dimensional Prefix Sums and Difference Arrays}
2D cumulative sums for $O(1)$ rectangle queries via inclusion-exclusion; difference arrays as the inverse operation, enabling $O(1)$ range updates.

\chapter{Kadane's Algorithm and Maximum-Subarray Variants}
Maximum subarray sum as a linear-time DP; variations (maximum product subarray, circular array, allowed deletions); the broader subarray-vs.-subsequence distinction.

\chapter{The Dutch National Flag Problem}
Three-way partitioning in a single pass with constant space, and its role inside three-way quicksort.

\chapter{Set Operations on Arrays}
Intersection and union via two-pointer (sorted) or hash-set (unsorted) techniques; symmetric difference; two-way and $k$-way (heap-based) merging of sorted arrays.

\chapter{Array Compaction and Element Removal}
Removing elements while preserving order (stable compaction); removing duplicates or specific elements in place.

\chapter{Element Frequency and the Majority Element}
Frequency counting and mode-finding; the Boyer-Moore majority-vote algorithm for finding a majority element in linear time and constant space.

\chapter{In-Place Rearrangement and Constant-Space Techniques}
Alternating positive/negative, even/odd segregation, and other custom-order rearrangements; the general toolbox of index mapping, swapping, and iterative in-place transformation.

% ============================================================
% STAGE 4: Alignment, Views, and Vectorized Operations
%
%   Goal check: after this stage, the reader understands the layer
%   between "an array sits in memory" (Stage 1) and "a SIMD/GPU unit
%   processes it" (Volume II, Stage 6) -- the views, alignment rules, and
%   bulk operations that make vectorized array code possible.
% ============================================================

\chapter{Memory Alignment for Cache and SIMD}
Aligning to 64-byte cache-line boundaries and false-sharing prevention; 16/32/64-byte SIMD alignment requirements (\texttt{aligned\_alloc}, \texttt{posix\_memalign}); compiler attributes and pragmas for enforcing alignment.

\chapter{Bounds Checking: Strategies, Performance, and Language Support}
Runtime and compiler-generated bounds checks; measuring and reducing their overhead; how Rust's borrow checker, C++ \texttt{std::array}, and managed languages approach safety differently.

\chapter{Array Views, Slices, and Strided Access}
Non-owning references and views vs. copies; slicing syntax and semantics (start:stop:step, negative indices); strided access and the vectorization challenges it introduces; windows and subviews without copying.

\chapter{Array Broadcasting}
Implicit dimension expansion and shape-compatibility rules (NumPy-style); virtual expansion that avoids physically duplicating memory.

\chapter{Reduction and Scan Operations}
Sum, product, min/max, and custom reductions; inclusive scan (cumulative sum) and exclusive scan as the related but distinct primitive.

\chapter{Gather and Scatter Operations}
Indexed reads (gather) and indexed writes (scatter) via an index array, including conflict handling on scatter; hardware gather/scatter support in AVX2/AVX-512.

\chapter{Array Comprehensions, Generators, and Iterators}
Declarative array construction (list comprehensions, functional style) versus lazy generators and the iterator protocol, and the memory-efficiency tradeoff between them.

\chapter{Reshaping and Transposition}
Changing dimensions without copying data (views, Fortran vs. C order); in-place transpose for square arrays and out-of-place transpose for rectangular ones.

\chapter{Cache-Oblivious Transpose}
Recursive divide-and-conquer transposition that achieves optimal cache usage with no tuning parameters.

\chapter{Concatenation, Splitting, Tiling, and Padding}
Joining arrays along an axis; dividing into equal or unequal chunks; blocked/tiled decomposition for cache optimization; border, alignment, and ghost-cell padding.

\chapter{Specialized Matrix Structures}
Toeplitz and circulant matrices (fast multiplication via FFT); triangular and symmetric matrices (packed storage, BLAS conventions); band matrices and sparse-diagonal storage for tridiagonal systems.

\chapter{Vectorization Fundamentals and Loop Vectorization}
Writing vector-friendly code and enabling autovectorization; data dependencies, trip count, alignment, and remainder loops that block or enable it.

\chapter{Lazy Evaluation and Array Expressions}
Expression templates and deferred computation (as in the Eigen library); fusing chained array operations instead of materializing intermediates.

\chapter{Versioning, Copy-on-Write, and Immutable Arrays}
Persistent arrays that share data between versions; structural sharing and the benefits of immutability in functional-style array code.

\chapter{Performance Profiling, the Roofline Model, and Benchmarking}
Measuring access patterns and using hardware counters; arithmetic intensity and the roofline model's bandwidth/compute bounds; sound microbenchmarking methodology (warming, timing, statistical significance).

% ============================================================
% STAGE 5: Foundational Structures Built on Arrays
%
%   Goal check: this is the core of the volume. After this stage, the
%   reader can implement and reason about the classic data structures and
%   algorithms (stacks, queues, heaps, hash tables, strings, tries,
%   graphs, segment/Fenwick trees, dynamic programming, greedy,
%   backtracking) that are themselves built on top of arrays. This stage
%   keeps the most chapters because it IS the goal, not a side path to it.
% ============================================================

\chapter{Stacks: Implementation and Applications}
The general ADT notion of interface vs. implementation; LIFO push/pop on an array with dynamic resizing; applications to expression evaluation, parenthesis matching, and the function call stack.

\chapter{Monotonic Stacks}
Maintaining a monotonic property for amortized $O(n)$ "next greater element" queries; applications to the largest-rectangle-in-histogram and stock-span problems.

\chapter{Queues and Circular Buffers}
FIFO via a circular array with front/rear pointers and full-vs.-empty detection; applications to BFS, task scheduling, and producer-consumer buffering.

\chapter{Deques and Monotonic Queues}
Double-ended queues with circular-array implementation; the monotonic queue (deque-based) for sliding-window maximum/minimum in amortized $O(n)$.

\chapter{Priority Queues and Binary Heaps}
Priority-based insert/find-min/delete-min; the complete binary tree stored in an array with parent/child index formulas; insert (sift-up), extract-min (sift-down), decrease-key, and $O(n)$ build-heap.

\chapter{Heap Variants and Priority-Queue Applications}
$d$-ary heaps and their cache-performance tradeoff against binary heaps; priority queues in Dijkstra's algorithm, Prim's MST, event simulation, and job scheduling.

\chapter{Hash Tables: Fundamentals and Hash Functions}
Hash function, collision resolution, load factor, and dynamic resizing as the core design space; division, multiplication, and universal hashing.

\chapter{Collision Resolution: Chaining and Open Addressing}
Separate chaining with linked or dynamic-array buckets; linear, quadratic, and double hashing under open addressing, and the clustering problems each causes; advanced schemes (Robin Hood, hopscotch, cuckoo hashing).

\chapter{Hash Table Resizing and Performance Analysis}
Doubling and rehashing with amortized and incremental strategies; how load factor governs expected chain length, probe sequence length, and worst-case behavior.

\chapter{Perfect and Minimal Perfect Hashing}
Collision-free hashing for static key sets via two-level construction; space-optimal minimal perfect hash functions (e.g. the CHD algorithm).

\chapter{Bloom Filters}
Probabilistic set membership via a bit array and $k$ hash functions; choosing optimal $k$ and size for a target false-positive rate; counting, scalable, and cuckoo-filter variants.

\chapter{Cardinality and Frequency Sketches}
Count-Min sketch for frequency estimation; HyperLogLog for approximate distinct counting in logarithmic space; comparing these and Bloom filters as a family of space-accuracy tradeoffs.

\chapter{Strings as Arrays and Manipulation Primitives}
Null-terminated vs. length-prefixed representations; concatenation, substring, and character-level access and modification as array operations.

\chapter{Naive and Knuth-Morris-Pratt String Matching}
Brute-force matching and its $O(nm)$ cost; KMP's failure function and how it eliminates redundant comparisons for $O(n+m)$ matching.

\chapter{Boyer-Moore and Its Variants}
The bad-character and good-suffix rules with right-to-left scanning for sublinear average-case matching; the simplified, practically faster Boyer-Moore-Horspool variant.

\chapter{Rabin-Karp and Aho-Corasick}
Rolling-hash fingerprint matching, including multiple-pattern matching with collision handling; the Aho-Corasick automaton (a trie with failure links) for matching many patterns at once.

\chapter{Suffix Arrays and the LCP Array}
Sorting suffixes via $O(n \log n)$ and linear-time (SA-IS) construction; pattern matching and longest-common-substring queries; the LCP array and Kasai's algorithm.

\chapter{Suffix Trees vs. Suffix Arrays}
The space-time tradeoffs between the two representations and when each is the practical choice.

\chapter{Tries and Radix Trees}
Array-of-children trie nodes and their space overhead; path-compressed tries (Patricia/radix trees) for space efficiency; applications to autocomplete, spell checking, and IP routing.

\chapter{Ternary Search Trees}
A space-efficient trie variant with BST-like node structure and balanced performance between tries and hash tables.

\chapter{Graph Representations: Matrix, List, and Edge List}
Adjacency matrices ($O(V^2)$, $O(1)$ edge checks, dense graphs) versus adjacency lists ($O(V+E)$, sparse graphs) versus edge lists; implementation variants and matrix-based operations like transitive closure.

\chapter{Graph Traversal: BFS and DFS}
Queue-based breadth-first traversal for level-order exploration and shortest paths in unweighted graphs; stack-based (explicit or recursive) depth-first traversal for cycle detection and ordering.

\chapter{Topological Sorting}
DFS-based and Kahn's (BFS-based) algorithms, and their use in dependency resolution.

\chapter{Single-Source Shortest Paths: Dijkstra and Bellman-Ford}
Dijkstra's priority-queue-based relaxation for non-negative weights; Bellman-Ford's handling of negative weights and negative-cycle detection.

\chapter{All-Pairs Shortest Paths: Floyd-Warshall}
Dynamic programming directly on the adjacency matrix, and why $O(V^3)$ is the right tradeoff for dense graphs.

\chapter{Minimum Spanning Trees: Prim's and Kruskal's Algorithms}
Prim's priority-queue-driven tree growth versus Kruskal's edge-sorting-and-union-find approach, and when each is preferable.

\chapter{Disjoint-Set Union (Union-Find)}
Array-based parent pointers with path compression and union by rank/size; the (near-constant, inverse-Ackermann) amortized analysis; applications to Kruskal's MST and dynamic connectivity.

\chapter{Segment Trees}
A binary tree over an array for range queries; $O(n)$ construction, $O(\log n)$ range query and point update; lazy propagation for range updates, and variants (RMQ, range sum, arbitrary associative operations).

\chapter{Fenwick Trees (Binary Indexed Trees)}
Cumulative-frequency queries and updates in $O(\log n)$ using bit manipulation on indices; the dual representation that supports range updates with point queries.

\chapter{Sparse Tables and Range Query Structures Compared}
Preprocessing for $O(1)$ range-minimum queries on immutable arrays; comparing segment trees, Fenwick trees, and sparse tables on mutability, space, and query time.

\chapter{Dynamic Programming: Foundations and State Representation}
Optimal substructure and overlapping subproblems; memoization vs. tabulation; 1D, 2D, and multi-dimensional state spaces and rolling-array space optimization.

\chapter{Knapsack Problems}
0/1 knapsack via a 2D DP array; the unbounded variant collapsing to 1D; subset sum as a Boolean DP and its space-optimized form.

\chapter{Longest Common and Longest Increasing Subsequences}
LCS via a 2D DP table with traceback reconstruction; LIS via both the $O(n^2)$ DP and the $O(n \log n)$ patience-sorting approach.

\chapter{Edit Distance}
The Levenshtein 2D DP over insertions, deletions, and substitutions; Hirschberg's linear-space divide-and-conquer reconstruction.

\chapter{Matrix Chain Multiplication and Optimal Binary Search Trees}
Optimal parenthesization via 2D DP; weighted search cost minimization for optimal BSTs, including Knuth's optimization.

\chapter{Palindrome and Grid DP Patterns}
Longest palindromic substring and palindrome partitioning; minimum path sum, unique paths, and the broader family of grid-traversal DP problems.

\chapter{DP with Bitmasks}
Representing subsets as bits for problems like the traveling salesman, Hamiltonian path, and assignment problem.

\chapter{Divide and Conquer: Recurrences and Classic Analyses}
Recurrence relations and the master theorem; revisiting merge sort, quicksort, and binary search (Stage 2) through the recurrence-analysis lens.

\chapter{Closest Pair and Convex Hull}
Divide-and-conquer closest-pair in $O(n \log n)$; Graham scan, Jarvis march, and quickhull for convex hull, built on point-array and orientation-test primitives.

\chapter{Greedy Algorithms: Theory and Activity Selection}
The greedy-choice property, optimal substructure, and exchange-argument proofs; interval scheduling via earliest finish time as the canonical example.

\chapter{Fractional Knapsack and Huffman Coding}
Greedy-by-ratio fractional knapsack; Huffman's optimal prefix-free codes built via a priority queue.

\chapter{Backtracking Fundamentals and N-Queens}
Exhaustive search with pruning over a state-space tree; the N-Queens problem as the canonical backtracking example over an array representation.

\chapter{Subset and Permutation Generation}
Generating all subsets (backtracking, bit manipulation) and all permutations (backtracking, Heap's algorithm).

\chapter{Constraint Satisfaction: The Sudoku Solver}
Backtracking with constraint propagation over an array-based board representation.

\chapter{Online Algorithms and Competitive Analysis}
Online vs. offline algorithms and the competitive ratio; cache replacement (LRU, FIFO, optimal) as a worked example.

\chapter{Streaming Algorithms: Heavy Hitters and Distinct Elements}
One-pass, sublinear-space algorithms; finding frequent elements (Misra-Gries), and approximate distinct counting (Flajolet-Martin).

\chapter{Reservoir Sampling}
Maintaining a uniform random sample from a single-pass stream of unknown size.

\chapter{Amortized Analysis: Aggregate, Accounting, and Potential Methods}
The three standard proof techniques for amortized cost, illustrated on the dynamic-array growth example from Stage 1 and on a binary counter.

\chapter{Cache-Oblivious Algorithms}
Achieving optimal cache usage without knowing cache parameters (the ideal-cache model); recursive matrix transpose and cache-oblivious searching/sorting as worked examples.

% ============================================================
% STAGE 6: Parallel and Distributed Array Processing
%
%   Goal check: after this stage, the reader can take an array algorithm
%   from Stages 1-5 and scale it correctly across cores, GPUs, and
%   machines -- and knows the concurrency-correctness pitfalls along the
%   way. The real-time and energy clusters are kept brief: relevant to
%   arrays, but not the volume's center of gravity.
% ============================================================

\chapter{Concurrency, Parallelism, and Performance Metrics}
Concurrent vs. parallel execution; speedup, efficiency, and the work/span (DAG) model of parallelism.

\chapter{Amdahl's Law and Gustafson's Law}
The serial fraction's limit on speedup versus scaled speedup under growing problem size; strong vs. weak scaling.

\chapter{Processes, Threads, and Thread Pools}
Separate address spaces vs. shared memory; thread creation and joining (pthreads, \texttt{std::thread}); worker pools and task queues for amortizing creation overhead.

\chapter{Work Stealing and Fork-Join Parallelism}
Dynamic load balancing via stealing from other threads' deques (the Cilk model); recursive divide-and-conquer parallelism with fork/join synchronization.

\chapter{OpenMP: Loops, Data Sharing, and Reductions}
The parallel-for directive and scheduling policies (static, dynamic, guided); shared/private/firstprivate/lastprivate clauses; reduction operations and their private-copy combination.

\chapter{OpenMP: Tasks, SIMD, and the Memory Model}
Task-based irregular parallelism (\texttt{task}, \texttt{taskwait}, \texttt{taskloop}); SIMD loop directives; flush operations and thread-safe memory ordering.

\chapter{Data Races and Memory Consistency Models}
Concurrent read-write and write-write conflicts; sequential consistency, TSO, and weak/relaxed models; the happens-before relation and C++11's memory-order options.

\chapter{Locks: Mutexes, Spinlocks, and Reader-Writer Locks}
Mutual exclusion and deadlock potential; busy-waiting spinlocks vs. blocking mutexes and when each wins; reader-writer locks for read-heavy workloads.

\chapter{Semaphores, Barriers, and Condition Variables}
Counting and binary semaphores for signaling; synchronization barriers; the wait/notify pattern for avoiding busy-waiting in producer-consumer queues.

\chapter{Lock-Free Programming and Atomic Operations}
Compare-and-swap, the ABA problem, and the wait-free/lock-free distinction; atomic fetch-and-add and compare-exchange with their memory-ordering implications.

\chapter{Lock-Free Stacks and Queues}
CAS-based push/pop and the Michael-Scott queue; hazard pointers and epoch-based reclamation for safely freeing memory under concurrent access.

\chapter{Transactional Memory: Software and Hardware}
Atomic blocks with optimistic concurrency and rollback; software (versioning, validation, commit) versus hardware (Intel TSX, capacity limits and fallback) implementations.

\chapter{False Sharing}
Cache-line contention between logically independent variables, its performance cost, and how to detect and eliminate it.

\chapter{GPU Array Algorithms: Sorting, Reduction, and Scan}
Bitonic and radix sort on GPU; tree-based reduction with warp-level shuffle primitives; the work-efficient Blelloch scan and bank-conflict avoidance.

\chapter{GPU Matrix and Stencil Operations}
Tiled matrix multiplication with shared memory (cuBLAS); halo exchange and shared-memory tiling for stencil computations.

\chapter{Heterogeneous Computing and OpenCL}
CPU-GPU task partitioning and host-device coordination; OpenCL's platform/device/context/queue model and its global/local/private/constant memory spaces.

\chapter{Distributed Computing Fundamentals and MPI Basics}
Distributed memory, message passing, latency and bandwidth; MPI communicators, ranks, and basic send/receive.

\chapter{MPI: Point-to-Point and Non-Blocking Communication}
Blocking vs. non-blocking send/receive (\texttt{MPI\_Isend}/\texttt{MPI\_Irecv}); overlapping communication with computation.

\chapter{MPI: Collective Operations and Derived Datatypes}
Broadcast, scatter, gather, all-to-all, and reduction; \texttt{MPI\_Reduce}/\texttt{MPI\_Allreduce}; describing complex layouts with derived datatypes.

\chapter{Distributed Array Partitioning and Domain Decomposition}
Block, cyclic, and block-cyclic distribution; 1D/2D/3D domain decomposition and minimizing the surface-to-volume communication ratio.

\chapter{Ghost Cells, Halo Exchange, and Load Balancing}
Boundary-region communication for stencil computations; static vs. dynamic load balancing, including work stealing across nodes for irregular workloads.

\chapter{Parallel I/O, Checkpointing, and Fault Tolerance}
Collective I/O and parallel file systems; periodic checkpoint/restart; replication and algorithm-based fault tolerance.

\chapter{MapReduce and Hadoop}
The map/shuffle/reduce model for data-parallel array processing at scale; HDFS block storage, replication, and data locality.

\chapter{Spark, RDDs, and Distributed Array Operations}
Resilient distributed datasets and in-memory transformations/actions; DataFrame array columns and broadcast variables.

\chapter{Partitioned Global Address Space (PGAS) Languages}
A global logical view with local control over data placement (UPC, Chapel, X10); distributed array declarations and affinity.

\chapter{Distributed Consensus and Eventual Consistency}
Paxos and Raft for agreement under failure; CRDTs and conflict resolution as an alternative for convergent distributed arrays.

\chapter{Real-Time Systems and Scheduling}
Hard vs. soft deadlines and worst-case execution time (WCET) analysis for array operations; rate-monotonic and earliest-deadline-first scheduling; cache locking and scratchpad memory for predictable access.

\chapter{Energy-Efficient and Green Computing}
Dynamic voltage and frequency scaling and power-aware algorithm design; minimizing data movement for energy savings; energy proportionality as a system-level goal.

\chapter{Parallel Debugging and Performance Analysis Tools}
Race detection (ThreadSanitizer, Valgrind) and deterministic replay; profiling parallel programs (perf, VTune, HPCToolkit).

\chapter{Scalability Analysis and Parallel Design Patterns}
Strong vs. weak scaling plots and identifying bottlenecks; the communication-to-computation ratio; recurring design patterns (task parallelism, data parallelism, pipeline, master-worker, SPMD).

% ============================================================
% STAGE 7: Arrays in Science, Machine Learning, and Beyond
%
%   Goal check: this stage is explicitly a tour, not a textbook on any one
%   of these fields. Each chapter answers "how does this domain lean on
%   array representation and the techniques from Stages 1-6," not "teach
%   me this entire field." Compressed hardest here, per the chosen trim
%   level: the same domains as the original draft are still represented,
%   but each gets one consolidated chapter instead of a multi-chapter
%   deep dive that competed with the volume's actual subject.
% ============================================================

\chapter{Tensors and Tensor Computing Frameworks}
Tensors as a generalization of arrays (rank, Einstein summation); NumPy, TensorFlow, PyTorch, and JAX as computational-graph frameworks built on this generalization. (Broadcasting itself was already covered in Stage 4.)

\chapter{Neural Networks as Array Transformations}
Weight matrices, bias vectors, and layer operations as array transforms; convolution (direct, im2col, FFT-based), pooling, batch normalization, and dropout as the standard building blocks of a CNN.

\chapter{Backpropagation and Automatic Differentiation}
The chain rule as a backward pass over the computational graph; forward- and reverse-mode automatic differentiation; gradient arrays and optimizers (SGD, momentum, Adam) as array updates.

\chapter{Training at Scale: Batching, Pipelines, and Parallelism}
Mini-batching and its memory/parallelism tradeoffs; data-loading pipelines and prefetching; model parallelism and data parallelism (gradient synchronization, all-reduce) for distributed training.

\chapter{Model Compression: Quantization, Pruning, and Distillation}
Post-training and quantization-aware training (INT8, mixed precision); structured and unstructured weight pruning; teacher-student knowledge distillation.

\chapter{Attention and Transformer Architectures}
Query-key-value arrays and scaled dot-product attention; the transformer's self-attention and positional encoding; sparse/local attention variants that reduce the quadratic cost.

\chapter{Tensor Cores and Hardware Acceleration for ML}
Mixed-precision matrix multiplication on NVIDIA Tensor Cores and TPU systolic arrays, and how they map onto the array operations from earlier chapters.

\chapter{Matrix Multiplication Optimization and BLAS/LAPACK}
Blocking/tiling and Strassen's algorithm for matrix multiply; the BLAS Level 1/2/3 hierarchy and LAPACK's decomposition routines as the standard interface vendors optimize against.

\chapter{Dense and Sparse Linear Algebra}
Matrix decompositions (LU, QR, Cholesky, SVD) and numerical stability; sparse matrix-vector multiply and iterative solvers (Jacobi, conjugate gradient, GMRES) with preconditioning.

\chapter{Scientific Computing: Finite Differences and Finite Elements}
Discretizing derivatives via stencil operations (explicit vs. implicit, stability); mesh-based finite element assembly into sparse matrix systems.

\chapter{Computational Fluid Dynamics and N-Body Simulations}
Grid-based CFD simulation at scale; particle-array force calculations and the Barnes-Hut and fast-multipole approximations for $N$-body problems.

\chapter{The Fast Fourier Transform and Its Applications}
The Cooley-Tukey algorithm and bit-reversal permutation; convolution via FFT and 2D/3D FFTs for image and signal processing.

\chapter{Image Processing Fundamentals and Filtering}
Images as 2D/3D pixel arrays; convolution-based filtering (blur, edge detection), morphological operations, geometric transformations, and keypoint/feature detection.

\chapter{Video and Audio Signal Processing}
Frame arrays and motion estimation for video; waveform arrays, windowing, and the short-time Fourier transform for audio; FIR/IIR digital filters and spectral analysis.

\chapter{Time Series Analysis and Forecasting}
Trend and seasonal decomposition of temporal arrays; moving averages, exponential smoothing, and wavelet transforms as a multiresolution alternative.

\chapter{Geospatial and Climate Data Arrays}
Raster GIS data and spatial indexing; satellite/multispectral imagery; grid-based climate and weather simulation as a large-scale array workload.

\chapter{Arrays in Bioinformatics: Sequences and Alignment}
DNA sequences as character arrays and $k$-mers; De Bruijn-graph genome assembly; dynamic-programming sequence alignment (Smith-Waterman, BLAST).

\chapter{Arrays in Bioinformatics: Expression and Molecular Simulation}
Gene-expression matrices and differential expression analysis; mass-spectrometry peak arrays for proteomics; particle position/velocity arrays in molecular dynamics.

\chapter{Quantum Computing: Arrays of Complex Amplitudes}
Qubits and quantum gates as unitary-matrix operations on complex-valued state arrays; a brief tour of Grover's and Shor's algorithms and how statevector simulation maps to classical array hardware.

\chapter{DNA Data Storage}
Encoding binary data into nucleotide sequences with error correction; mapping bits to bases and the addressing challenge of random access.

\chapter{Neuromorphic and Analog Computing}
Spiking neural networks and event-driven computation; memristor crossbar arrays for in-memory analog matrix multiplication; processing-in-memory as a way to avoid the von Neumann data-movement bottleneck.

\chapter{Array Programming Languages: APL, J, and Julia}
The array-oriented paradigm's implicit parallelism and concise notation (APL/J rank polymorphism); Julia's multiple dispatch and broadcasting as a modern take on the same idea.

\chapter{NumPy Internals and Performance}
The \texttt{ndarray} structure, strides, and ufuncs; vectorization, avoiding Python-level loops, and JIT acceleration (numba) for closing the performance gap.

\chapter{Pandas, Dask, and Distributed DataFrames}
Columnar DataFrame storage built on array backing; Dask's lazy evaluation and task graphs for parallel and out-of-core array computing.

\chapter{Array Databases and Columnar Storage}
Native array query languages (SciDB, TileDB) with chunking and distribution; column-store layout and vectorized, SIMD-accelerated query execution.

\chapter{JIT Compilation for Array Workloads}
Runtime specialization (Numba, JAX's \texttt{jit}); XLA's fusion and optimization passes; polyhedral loop transformations for automatic parallelization.

\chapter{Memory-Mapped and Out-of-Core Arrays}
Using \texttt{mmap} and virtual memory to treat disk-backed arrays as in-memory ones; streaming and blocked processing for arrays larger than RAM.

\chapter{Persistent and Versioned Arrays}
Structural sharing and path copying for immutable, persistent arrays; copy-on-write versioning for git-like array history with efficient diffs.

\chapter{Compressed Array Encodings for Analytics}
Run-length, dictionary, delta, and frame-of-reference encoding as used in columnar analytical databases. (The general compression theory behind these techniques is covered in Volume I, Stage 17; this chapter focuses on the array/database-specific application.)

\chapter{Succinct Data Structures and Wavelet Trees}
Representing data in space close to the information-theoretic minimum while still supporting rank/select queries; the wavelet tree as a succinct structure for range queries and text indexing.

\chapter{The Van Emde Boas Layout and External-Memory Algorithms}
A recursive, cache-oblivious array layout for implicit search trees; the I/O model and algorithms (sorting, searching, graph algorithms) designed to minimize block transfers.

\chapter{Approximate Computing}
Deliberately trading accuracy for performance or energy under probabilistic guarantees and acceptable error bounds.

\chapter{Privacy-Preserving Array Computation}
Differential privacy (the Laplace and Gaussian mechanisms); homomorphic encryption and secure multi-party computation for computing on arrays without revealing them.

\chapter{Blockchain, Merkle Trees, and Distributed Ledgers}
Hash trees over arrays for efficient integrity proofs; append-only ledger arrays, consensus, and Byzantine fault tolerance.

\chapter{Edge Computing, Federated Learning, and TinyML}
Local array processing under resource constraints; federated training with gradient aggregation instead of centralized data; quantized neural networks running on microcontrollers.

\chapter{Hardware Accelerators: FPGAs and Beyond}
FPGAs, ASICs, and other domain-specific architectures for array workloads; reconfigurable logic, custom datapaths, and high-level synthesis.

\chapter{Alternative Computing Paradigms: Optical, Reversible, and Biological}
Photonic matrix multiplication at the speed of light; reversible logic and the Landauer limit on energy efficiency; DNA computing and cellular automata as array-like biological computation.

\chapter{Looking Ahead: Future Memory and Post-Moore Computing}
Persistent and storage-class memory; 3D stacking, chiplets, and domain-specific architectures as the industry's response to the end of simple transistor scaling (Volume I, Stage 4); open problems in array computing.

